% FORMAT THIS
\chapter{Part 0 : Introduction}

This tutorial has been engineered to be \emph{fun}.

If you work through it and, for some reason, do \emph{not} have a resounding \mbox{``I-Rule''} feeling afterwards, please send us an email and tell us how to improve it: \href{mailto:contact@moflon.org}{contact@moflon.org}

\begin{figure}[htp]
\begin{center}
	% RENAME THIS IMAGE ??
	\includegraphics[height=0.45\textheight]{../introduction_images/i-rule}
	\caption{How you should feel when you're done.}
	\label{i-rule}
\end{center}
\end{figure}
\break
 
 % End page ------------------------------------------------------------------------------------------------------------------------------------------------------------------------


To enjoy the experience, you should be fairly comfortable with Java or a comparable object-oriented language, and know how to perform basic tasks in Eclipse.  Although we assume this, we give references to help bring you up to speed as necessary.  
Last but not least, very basic knowledge of common UML notation would be helpful.

Our goal is to give a \emph{hands-on} introduction to metamodelling and graph transformations using our tool \emph{eMoflon}.
The idea is to \emph{learn by doing} and all concepts are introduced while working on a concrete example.
The language and style used throughout is intentionally relaxed and non-academic.
For those of you interested in further details and the mature formalism of graph transformations, we give relevant references throughout the tutorial.

This tutorial has been split into 6 parts. Below are a few suggestions of what you must read, or what you can skip depending on what you're interested in.
 
% WRITE BRIEF (for now) DESCRIPTIONS OF EACH PART + APPROXIMATELY HOW LONG IT SHOULD TAKE
\begin{description}

\item[Part I: Installation and Set Up]provides a very simple example and a few JUnit tests to test the installation and configuration of eMoflon.
 
After working through this chapter, you should have an installed and tested eMoflon working for a trivial example.
We also explain the general workflow and the different workspaces involved.

This chapter can be considered \emph{mandatory} if you are new to eMoflon and we recommend working through it in any case.
It's also kept as minimal as possible and should only take a few minutes really.

\item[Part II: Ecore]  Describe the Ecore section here; is the main chapter and takes you step-by-step through a more realistic example that showcases most of the features we currently support.

Working through this chapter should serve as a basic introduction to model-driven engineering, metamodelling and graph transformations.

You should definitely work through this chapter if you're new to metamodelling in general (using Ecore/EMF)

\item[Part III: SDMs] Describe the part here ; you should read this if you're interested in model transformation via graph transformation using Story Driven Modeling

\item[Part IV: TGGs] Describe the part here; Feel free to jump directly here after working through at least Part I if you're mainly interested in TGGs.
Although the example builds up on parts constructed in previous chapters, we provide a ``cheat'' package that you can use to get started directly.

\item[Part V: Model To Text Transformations] Describe the part here;

\item[Part VI: Miscellaenous] Describe the part here; Great references files to keep on hand while using ECore and it's partner programs. Help to avoid mistakes. (Grokking EA)

If you're in a hurry, this chapter can be skipped

\end{description}

One last thing: at the moment we unfortunately only support Windows. This should hopefully change in future releases.

Well, that's it -- sit back, relax, grab a coffee and enjoy the ride!


